Les tres gràfiques del dibuix són les d'una funció $f$, de la seva derivada $f'$ i d'una primitiva $F$
de $f$, en un ordre qualsevol.
\begin{center}
\foreach \nom/\funcio/\dom in {A/{\x*\x-1}/{-2.05:2.05},B/{2*\x}/{-1.6:1.6},C/{\x*\x*\x/3-\x}/{-2.25:2.25}}{%
\begin{tikzpicture}[x=0.55cm,y=0.45cm]
  \draw[gray!55,very thin,step=1] (-2.5,-3.5) grid (2.5,3.5);
  \draw[->] (-2.7,0) -- (2.9,0) node[below,font=\scriptsize] {$x$};
  \draw[->] (0,-3.7) -- (0,3.9) node[left,font=\scriptsize] {$y$};
  \foreach \i in {-2,-1,1,2} \draw (\i,0.12) -- (\i,-0.12) node[below,font=\tiny] {$\i$};
  \foreach \j in {-2,2} \draw (0.1,\j) -- (-0.1,\j) node[left,font=\tiny] {$\j$};
  \draw[\colorgrafica,very thick,domain=\dom,samples=60] plot (\x,\funcio);
  \node[font=\small\bfseries] at (0,-4.6) {\nom};
\end{tikzpicture}\hspace{0.8cm}}
\end{center}

\begin{apartats}

\apartat[1,25]{0,75}
Digues quina gràfica és $f$, quina és $f'$ i quina és $F$. Justifica-ho.

\begin{solucio}
On una funció té un extrem relatiu, la seva derivada s'anul·la. La C té extrems a $x=-1$ i $x=1$, on la
A val 0; i la A té un mínim a $x=0$, on la B val 0. Per tant, la A és la derivada de la C, i la B és la
derivada de la A: $F$ és la C, $f$ és la A i $f'$ és la B.\\
També es veu en el creixement: la C creix on la A és positiva, i la A creix on la B és positiva.
\end{solucio}

\begin{tria}{tres-grafiques}
\itemtria{expressio}{1}{1,25}
La gràfica A és la de $y=x^2-1$. Troba la primitiva de $x^2-1$ que passa per l'origen, i comprova que
té els extrems relatius on els té la gràfica C.

\begin{solucio}
$F(x)=\frac{x^3}{3}-x+k$, i $F(0)=0$ dona $k=0$: $F(x)=\frac{x^3}{3}-x$. $F'(x)=x^2-1=0$ a $x=\pm1$:
màxim $F(-1)=\frac23$ i mínim $F(1)=-\frac23$, els de la gràfica C.
\end{solucio}

\itemtria{altra-primitiva}{1}{1,25}
La gràfica A és la de $y=x^2-1$. Troba la primitiva $G$ de $x^2-1$ que passa pel punt $(0,2)$. Quina
relació té la seva gràfica amb la C? On té els extrems relatius?

\begin{solucio}
$G(x)=\frac{x^3}{3}-x+k$ i $G(0)=k=2$: $G(x)=\frac{x^3}{3}-x+2$. Com que $G$ i la C es diferencien en la
constant 2, la gràfica de $G$ és la C desplaçada 2 unitats cap amunt. Té els extrems a les mateixes
abscisses: màxim $\left(-1,\frac83\right)$ i mínim $\left(1,\frac43\right)$.
\end{solucio}
\end{tria}

\begin{nomesllarg}
\apartat{0,75}
Hi ha alguna primitiva de la funció A que tingui un sol extrem relatiu? Raona-ho.

\begin{solucio}
No. Totes les primitives de $f$ són $F(x)+k$, i tenen la mateixa derivada $f$, que s'anul·la i canvia
de signe a $x=-1$ i a $x=1$. Totes tenen un màxim a $x=-1$ i un mínim a $x=1$.
\end{solucio}
\end{nomesllarg}

\end{apartats}
